\emph{Frames.} A change of local frame acts on the matrix of $G$ by congruence, $G \mapsto P^{T} G P$ with $P$ invertible. Congruence is a diffeomorphism of the symmetric matrices that preserves corank and signature (\claimref{synthesis.completion-recovers-the-real-clifford-algebras}(iii)), so it preserves each corank stratum, and transversality to a stratum is frame-independent; $\det G$ changes by the nonvanishing factor $\det(P)^2$, so its zero set and the order of vanishing are frame-independent.

\emph{Step 1: the Schur normal form.} Let $x$ have corank $r$ and choose a frame with $G_x = 0_r \oplus B_0$, $B_0$ nondegenerate of signature $(p_0, q_0)$. Near $x$ write $G = \begin{pmatrix} A & C \\ C^{T} & B \end{pmatrix}$ with $B$ invertible. Congruence by $\begin{pmatrix} 1 & 0 \\ -B^{-1}C^{T} & 1 \end{pmatrix}$ carries $G$ to $S \oplus B$ with $S := A - C B^{-1} C^{T}$, a smooth $r \times r$ symmetric matrix function with $S(x) = 0$. Hence near $x$ the corank of $G$ is the corank of $S$, the signature of $G$ is $\mathrm{sig}(S) + (p_0, q_0)$, and $\det G = \det S \cdot \det B$. The map $(A, C, B) \mapsto S$ is a submersion near $G_x$ (its derivative in the $A$ direction is the identity), so near $G_x$ the corank-$r$ matrices are the preimage of $0$, a submanifold of codimension $\dim \mathrm{Sym}_r = r(r+1)/2$, and $G$ is transverse at $x$ exactly when $d(S \circ G)_x : T_xM \to \mathrm{Sym}_r$ is surjective.

\emph{Lemma A (corank one).} The regularity characterisation of tameness, (T1)-(T3) as the regular zero set of $\det G$, and the $r + 1$ signature count of Lemma B are also proved in the result filed alongside, synthesis.tameness-splits-regularity-and-order; they are repeated here in the form the stability argument uses. At corank one, $G_x\,\mathrm{adj}(G_x) = \det(G_x) = 0$, so every column of the adjugate lies in the kernel, spanned by $k$; the adjugate is symmetric and nonzero (a rank $n-1$ matrix has a nonzero $(n-1)$-minor), so $\mathrm{adj}(G_x) = c\,k k^{T}$ with $c \neq 0$. Jacobi's formula gives $d(\det G)_x = \mathrm{tr}(\mathrm{adj}(G_x)\, dG_x) = c\, k^{T} dG_x\, k$. Extending $k$ to a covector field and putting $g_k = G(k,k)$, $d(g_k)_x = k^{T} dG_x k + 2 k^{T} G_x\, dk = k^{T} dG_x k$, since $G_x k = 0$; this is $\delta(k \otimes k)$ at $x$ (Lemma \claimref{sheaf.invariant-form-t3}). So $d(\det G)_x = c\,\delta(k \otimes k)$, and (T3) of Hypotheses \claimref{sheaf.tameness} at $x$ is $d(\det G)_x \neq 0$. For $r = 1$ the stratum is locally $\{\det = 0\}$ with $d\det$ nonzero on it ($H \mapsto c\,k^{T} H k$), so transversality is again $d(\det \circ G)_x \neq 0$. Then $\Sig = (\det G)^{-1}(0)$ is a smooth hypersurface near $x$; corank at most one is an open condition, so $n - 1$ eigenvalues of $G$ stay away from zero near $x$ and the remaining one has the sign of $\det G$ divided by their product, which changes across $\Sig$ because the zero of $\det G$ is simple. The two sides therefore carry $(p_0 + 1, q_0)$ and $(p_0, q_0 + 1)$. The check \texttt{check\_corank\_one\_differential\_is\_delta} verifies $\mathrm{adj}(G_0) = c\,kk^{T}$, $c \neq 0$, and that the $t$-coefficient of $\det(G_0 + tH)$ is $c\,H(k,k)$, for six corank-one cometrics and four directions each, including directions with $H(k,k) = 0$; \texttt{check\_tame\_wall\_is\_two\_signature} verifies the signature jump.

\emph{Lemma B (signatures), which is clause (2).} If $G$ is transverse at $x$ then $S \circ G$ is a submersion at $x$, so by the submersion theorem it maps every neighbourhood of $x$ onto a neighbourhood of $0$ in $\mathrm{Sym}_r$. That neighbourhood contains $\varepsilon\,\mathrm{diag}(1, \dots, 1, -1, \dots, -1)$ with $j$ entries $-1$, for each $j = 0, \dots, r$ and small $\varepsilon > 0$; there $G$ has signature $(p_0 + r - j, q_0 + j)$. For $r = 2$ these are $(p_0 + 2, q_0)$, $(p_0 + 1, q_0 + 1)$, $(p_0, q_0 + 2)$. If $(1,3)$ is among them then $(p_0, q_0)$ is $(0,2)$ or $(1,1)$ (the first entry cannot be $p_0 + 2 = 1$ with $p_0 + q_0 = 2$ unless $p_0 = -1$), and in both cases $(2,2)$ is in the list. The check \texttt{check\_transverse\_corank\_two\_three\_signatures} computes the three signatures exactly for three blocks $B$ at three scales, and the count $r + 1$ for $r = 1, 2, 3$.

\emph{Clause (1), the cone.} Transversality at corank two needs a surjection $T_xM \to \mathrm{Sym}_2 \cong \R^3$, so $n \geq 3$. By the submersion theorem there are coordinates $(y_1, y_2, y_3, z)$, $z \in \R^{n-3}$, centred at $x$ with $S \circ G = \begin{pmatrix} y_1 & y_2 \\ y_2 & y_3\end{pmatrix}$. Then $\det G = (y_1 y_3 - y_2^2)\det B$ with $\det B \neq 0$, which vanishes to second order at $x$, and $\Sig$ is the cone $y_1 y_3 = y_2^2$ (a quadric of signature $(1,2)$: $y_1 y_3 - y_2^2 = \tfrac14(y_1 + y_3)^2 - \tfrac14(y_1 - y_3)^2 - y_2^2$) times $\R^{n-3}$. It contains the lines through $x$ in the directions $(1,1,1,0)$, $(1,-1,1,0)$, $(4,2,1,0)$ and every $z$ direction, which together span $T_xM$; a smooth hypersurface through $x$ contains only curves tangent to its $(n-1)$-dimensional tangent space there, so $\Sig$ is not a smooth hypersurface at $x$ and (T1) fails. The radical of $G_x$ has rank two, so (T2) fails. The check \texttt{check\_cone\_is\_not\_a\_hypersurface} verifies, in the model with $n = 4$, that the linear coefficient of $\det G(tv)$ vanishes in six directions while a quadratic one does not, that the four lines lie in $\Sig$ and span $\R^4$, and that the corank-two locus is the $z$-line.

\emph{Clause (1), stability.} Let $G'$ be a cometric with $\|G' - G\|_{C^1} < \eta$ on a closed coordinate ball $\bar U$ about $x$. For small $\eta$ the block $B'$ stays invertible on $\bar U$ and the Schur complement $S'$ of $G'$ (a smooth function of the matrix entries) satisfies $\|S' - S\|_{C^1(\bar U)} \leq K\eta$ for a constant $K$ fixed by $G$ on $\bar U$. Fix $z_0$ near $0$ and work on the ball $|y| \leq \rho$ inside $\bar U$, where $S(y, z_0) = y$ in the identification $\mathrm{Sym}_2 \cong \R^3$ above. The map $T(y) = y - S'(y, z_0) = (S - S')(y, z_0)$ has Lipschitz constant at most $K\eta$ and moves no point by more than $K\eta$; for $K\eta \leq \min(\tfrac12, \tfrac{\rho}{2})$ it maps the ball into itself and is a contraction, so it has a unique fixed point $y^{*}$, $|y^{*}| \leq 2K\eta$, where $S' = 0$, that is where $G'$ has corank two. There $d_y S' = 1 - d_yT$ is within $\tfrac12$ of the identity, hence invertible, so $G'$ is transverse at $(y^{*}, z_0)$. The check \texttt{check\_corank\_two\_point\_survives\_perturbation} exhibits the persisting corank-two point exactly, by solving $S' = 0$, for two constant perturbations at $\varepsilon = 10^{-1}, 10^{-2}, 10^{-3}$ (with $|y^{*}| \leq 10\varepsilon$) and for a position-dependent perturbation of the $A$ block at three values of $z$.

\emph{Clause (1), conclusion, and clause (3).} The corank-two points of a cometric transverse to the corank-two stratum are the preimage of that stratum, a submanifold of codimension three. Let $G$ be transverse at every point of $\Sig$. At a point of corank $r \geq 2$ Lemma B puts at least three signatures in every neighbourhood and the cone argument shows (T1) and (T2) fail there (for $r > 2$ the radical has rank $r$ and (T2) fails directly). At a point of corank at most one some neighbourhood meets at most two signatures by Lemma A, and if every point of $\Sig$ has corank one then Lemma A gives (T1) and (T3) everywhere on $\Sig$ ($\Sig$ is the zero set of $\det G$, hence closed, and a regular zero set, hence a smooth hypersurface), while the radical lines are the kernel of a smooth family of constant rank $n - 1$ on $\Sig$ (locally the span of a nonzero column of the adjugate), a smooth line subbundle, which is (T2). So the points where tameness fails are exactly the points near which three or more signatures meet. For $n = 2$ there is no surjection $T_xM \to \mathrm{Sym}_2$, so a transverse cometric on a surface has no corank-two point and is tame. For $n \geq 3$ the model of the cone argument is transverse at the origin and by the stability argument every cometric $C^1$-close to it on $\bar U$ has a corank-two point, so it is not tame: the non-tame cometrics contain a nonempty open set, and tameness is not generic.

\emph{Remark (argued, not part of the claim).} Every transverse corank-two point adjacent to this world's signature $(1,3)$ opens a $(2,2)$ region, a residue-$0$ signature with two times. The series excludes two-time signatures for the $r = 0$ tower by an imported instability of two-time dynamics. If that exclusion is extended to regions of one spacetime, transverse corank-two points next to $(1,3)$ are ruled out, and for degenerations transverse along $\Sig$ the clauses above then give tameness near every point of $\Sig$ adjacent to $(1,3)$. Whether the exclusion may be read that way is a decision about the premise, not a theorem here.
